\section{SF-E835. Every semistable torsion-free balanced sheaf reduces
to the checked hull
gate}\label{SF-E835:sf-e835.-every-semistable-torsion-free-balanced-sheaf-reduces-to-the-checked-hull-gate}

Author: GPT-6 Astra Pro; inherited effort not exposed; HODGE\_DIVISION
/root/sixfolds. Date: 2026-11-21.

\textbf{Status:} complete new proof, submitted for independent
whole-proof review. Its new analytic premise is SF-E834, which is
separately frozen and under review. This section does not promote that
premise before its review. The remaining steps are algebraic and are
proved in full below. The statement concerns actual semistable
torsion-free sheaves; it does not extend to unstable sheaves or
arbitrary perfect complexes.

\subsection{1. Statement}\label{SF-E835:statement}

Retain the genuine polarized RM abelian fourfold X and
theta=theta\_++theta\_- of SF-E834. Let E be any theta-slope-semistable
torsion-free coherent sheaf of positive rank R, and suppose

\[
 e^B\operatorname{ch}(E)
   =R\left(1-\frac q2\theta_+^2\right)
      \left(1-\frac q2\theta_-^2\right),
 \qquad B\in\operatorname{NS}(X)_{\mathbb Q},\quad q\in\mathbb Q_{>0}.
 \tag{E835.1}
\]

Then the actual hull \(H=E^{\vee\vee}\) is locally free and has a
filtration by subbundles whose quotients are stable projectively flat
vector bundles of the same theta-slope. The quotient H/E is supported on
finitely many points or is zero. Consequently E is not weakly
G-semiregular for any specified coherent effective finite
translation/tensor action on E in the scope of SF-E818.

This supplies the hypotheses on the hull which were explicitly assumed
in E818. In particular a singular reflexive hull or a singular stable
factor in a locally free hull is no longer an unchecked semistable
escape, once E834 is established. No chosen stable-factor filtration is
assumed to be preserved by G.

\subsection{2. Stable torsion-free factors and stability of their
hulls}\label{SF-E835:stable-torsion-free-factors-and-stability-of-their-hulls}

Choose a finite slope-stable-factor filtration

\[
 0=E_0\subset E_1\subset\cdots\subset E_s=E,
 \qquad Q_i=E_i/E_{i-1},\qquad r_i=\operatorname{rk}Q_i>0,
 \tag{E835.2}
\]

with torsion-free quotients of slope mu\_theta(E). Here is a direct
existence argument adequate for this use. If a semistable torsion-free
sheaf is not stable, choose an equal-slope proper subsheaf of smallest
positive rank and saturate it. Saturation adds a torsion quotient whose
codimension-one cycle is effective. Any increase of slope would
contradict semistability of the ambient sheaf; thus the saturated sheaf
still has equal slope. Minimality of its rank makes it stable. The
saturated quotient is torsion-free and semistable: a larger-slope
subsheaf of the quotient would lift to a larger-slope subsheaf of the
ambient sheaf. Induction on rank gives(E835.2). This does not assume
local freeness of its factors.

Set

\[
 H_i=Q_i^{**},\qquad T_i=H_i/Q_i,
 \quad \alpha_i=\frac{c_1(Q_i)}{r_i}=\frac{c_1(H_i)}{r_i},
 \quad \alpha=\frac{c_1(E)}R,\quad\delta_i=\alpha_i-\alpha.
 \tag{E835.3}
\]

Every H\_i is stable reflexive. Indeed, for any proper positive-rank
subsheaf M of H\_i with 0\textless rank(M)\textless r\_i, the
intersection M with Q\_i has the same rank and first Chern class as M:
the quotient is supported in codimension at least two, since H\_i/Q\_i
is. Stability of Q\_i then gives mu(M)=mu(M intersection
Q\_i)\textless mu(Q\_i)=mu(H\_i). Each T\_i has support dimension at
most two.

The first Chern component of(E835.1) gives B=-alpha. Equality of the
factor slopes implies integral\_X delta\_i theta³=0. Hence each
invariant representative delta\_i is primitive for theta.

\subsection{3. Mixed discriminant equality eliminates all surface
defects}\label{SF-E835:mixed-discriminant-equality-eliminates-all-surface-defects}

Let Delta be the discriminant convention(E834.1). Additivity of Chern
characters in(E835.2) gives the exact variance identity

\[
 \Delta(E)=\sum_i\Delta(Q_i)-\frac12\sum_i r_i\delta_i^2
          =\sum_i\Delta(H_i)+\sum_i\operatorname{ch}_2(T_i)
                         -\frac12\sum_i r_i\delta_i^2.
 \tag{E835.4}
\]

For example, subtracting the first expression from the definition of
Delta(E) reduces to sum\_i r\_i(alpha\_i-alpha)²=sum\_i r\_i alpha\_i²-R
alpha², using sum\_i r\_i alpha\_i=R alpha. The second expression uses
c1(Q\_i)=c1(H\_i) and ch2(Q\_i)=ch2(H\_i)-ch2(T\_i).

Every mixed pairing on the right side is nonnegative:

\begin{itemize}
\item
  E834 applies to the stable reflexive H\_i.
\item
  The leading-character/support-cycle identity makes ch2(T\_i) the
  effective surface cycle with positive generic lengths. E818§1 proves
  that every nonzero effective surface has strictly positive pairing
  with theta\_+theta\_- on a genuine RM fourfold.
\item
  The elementary primitive mixed sum-of-squares formula E811.4 gives
  -integral\_X delta\_i² theta\_+theta\_- nonnegative.
\end{itemize}

The left side has zero mixed pairing, by(E835.1). Each nonnegative
summand therefore vanishes. E834's equality case makes every H\_i a
locally free projectively flat stable bundle. Strict mixed positivity
eliminates every two-dimensional component of T\_i. The equality case of
E811.4 makes every delta\_i block diagonal for V\_+ direct sum V\_- and
primitive on each plane separately. Thus

\[
 \operatorname{ch}(H_i)=r_i e^{\alpha_i},\qquad
 \dim\operatorname{Supp}T_i\le1,
 \qquad \delta_i=\delta_{i,+}+\delta_{i,-},\quad
 \delta_{i,\pm}\theta_\pm=0.
 \tag{E835.5}
\]

The forms in the last identities are invariant forms on the indicated
planes. No claim that those planes are rational subvarieties is made.
The complete strict-surface proof is E818§1, including its Pluecker and
Albanese argument, not a supposition of general-position positivity.

\subsection{4. The cubic character eliminates all curve
defects}\label{SF-E835:the-cubic-character-eliminates-all-curve-defects}

On each rank-two complex plane, delta\_i,plus³=delta\_i,minus³=0.
Therefore

\[
 \delta_i^3=3\delta_{i,+}^2\delta_{i,-}
                 +3\delta_{i,+}\delta_{i,-}^2,
 \qquad \theta\delta_i^3=0.
 \tag{E835.6}
\]

Multiplication by theta kills the two summands either by excess plane
degree or by the separate primitivity in(E835.5). The actual K0 identity
E=sum\_i H\_i-sum\_i T\_i and B=-alpha give

\[
 0=\left(e^{-\alpha}\operatorname{ch}(E)\right)_3
     =\frac16\sum_i r_i\delta_i^3-\sum_i\operatorname{ch}_3(T_i).
 \tag{E835.7}
\]

Here e\^{}-alpha leaves the third character of each T\_i unchanged
because T\_i already has support dimension at most one. Multiply by
theta and integrate. The first term is zero by(E835.6). Each ch3(T\_i)
is its effective curve cycle, with positive generic-length weights. The
ample theta has strictly positive degree on every nonzero effective
curve. Thus every T\_i has support dimension zero or is zero. Let P be
the finite union of their supports and put U=X minus P.

The exact leading-character source is the same as in E818: Fulton,
\emph{Intersection Theory}, second edition, Example18.3.11 p363 together
with Theorem18.3(3) p353, since the tangent bundle of X is trivial. No
sign is lost when replacing a sheaf by its effective leading support
cycle. The actual short exact sequences Q\_i into H\_i supply the
positive weights; this argument would not apply to virtual defect
differences.

\subsection{5. Extensions of vector bundles extend across these
points}\label{SF-E835:extensions-of-vector-bundles-extend-across-these-points}

We give the required algebraic extension step rather than assuming that
a reflexive hull preserves exact sequences. For any two vector bundles
V,W on the smooth projective fourfold X, restriction gives an
isomorphism

\[
 \operatorname{Ext}^1_X(V,W)
       \xrightarrow{\ \sim\ }
 \operatorname{Ext}^1_U(V|_U,W|_U).
 \tag{E835.8}
\]

Indeed both groups are H1 of the corresponding bundle V dual tensor W,
since the first argument is locally free. The reduced finite set P is
locally cut out by four regular parameters. For a locally free sheaf K,
its local-cohomology sheaves with support P vanish in degrees below4.
Explicitly, the augmented Cech complex on those four parameters computes
local cohomology; regularity of the sequence makes its cohomology vanish
except in degree4, and tensoring with the free stalk of K preserves this
statement. This is precisely Stacks Lemma36.6.9, tag0G7P, whose complete
proof was read. The local-to-global support spectral sequence therefore
gives H\_P\^{}j(X,K)=0 for j\textless4. The support triangle and its
long exact sequence, Stacks §20.21 tag0A39, now give H1(X,K) isomorphic
to H1(U,K\textbar U). This proves(E835.8). No analytic/algebraic
comparison on the nonproper U is used: the sheaves and the actual
extensions being restricted are algebraic; holomorphic sheaves on the
original projective X may first be algebraized by GAGA.

On U, each Q\_i agrees with H\_i. The restricted filtration(E835.2) is
therefore a filtration by vector bundles with vector-bundle quotients.
Inductively construct vector bundles K\_i on X extending E\_i\textbar U
with a filtration of the prescribed kind. Set K\_1=H\_1. Given K\_(i-1),
the actual sequence on U is

\[
 0\longrightarrow K_{i-1}|_U\longrightarrow E_i|_U
             \longrightarrow H_i|_U\longrightarrow0.
 \tag{E835.9}
\]

Its extension class lifts by(E835.8) to an actual sequence on X

\[
 0\longrightarrow K_{i-1}\longrightarrow K_i
                    \longrightarrow H_i\longrightarrow0.
 \tag{E835.10}
\]

The middle term is locally free: on an affine local ring the free
quotient H\_i splits the sequence. The new inclusion is a subbundle
inclusion, again because its quotient is locally free. Restriction to U
identifies(E835.10) with(E835.9), including the indicated subobject and
quotient. Induction produces K\_s with stable locally free quotients
H\_i of the same theta-slope.

The original hull H and K\_s are reflexive and agree on U. The
complement has codimension4, so the usual reflexive Hartogs extension
property identifies them uniquely. Thus H is K\_s and has the claimed
actual subbundle filtration. The canonical injection E into H is an
isomorphism on U, so H/E is supported on the finite set P.

\subsection{6. Applying the existing categorical
obstruction}\label{SF-E835:applying-the-existing-categorical-obstruction}

All hypotheses of SF-E818 have now been proved for E: its actual hull is
locally free with stable locally free same-slope factors, and it has the
same balanced character. If H/E is nonzero, the canonical double-dual
sequence and the given coherent action give the point-defect kernel of
SF-E817, of dimension at least1+3s for s support orbits. If H/E is zero,
SF-E813 excludes the locally free sheaf with this filtration.
Equivalently one may invoke the complete theorem E818 directly.

Neither the chosen stable-factor filtration(E835.2) nor its hull factors
need carry separate G actions. The only action used at the final step is
the specified coherent action on E and its canonical functorial actions
on H and H/E, as proved in E817/E818. The scalar twisting and Serre
adjunction inputs are inherited in exactly their reviewed scopes.

\subsection{7. Source and review
boundary}\label{SF-E835:source-and-review-boundary}

The new analytic premise E834 is explicitly separate and must pass its
own whole-proof reviews. The strict surface positivity, leading support
cycle, point-defect and locally free filtration results are the reviewed
E818/E817/E813 statements. All stable-hull, variance, curve, extension-
lifting and actual-filtration steps above are new complete arguments.
The exact local-cohomology inputs were freshly read in full at
\url{https://stacks.math.columbia.edu/tag/0G7P} and
\url{https://stacks.math.columbia.edu/tag/0A39}; the depth form was also
checked at tag0AVZ. No finite numerical certificate is needed or
claimed. An unstable sheaf or a perfect complex outside these actual
sheaf hypotheses remains an open construction route.
